% Solutions of 27-combinational-circuits.
\section*{Chapter \ref{ch:b1:combinational-circuits} --- Boolean Algebra and Combinational Circuits}

\begin{solution}{exo:b1:combinational-circuits:1}
$f = ab' + bc$ is 1 on the rows $abc = 011, 100, 101, 111$ and 0 on the others. The longest path is NOT, AND, OR: delay 3.
\end{solution}

\begin{solution}{exo:b1:combinational-circuits:2}
$f = a'b'c + a'bc' + ab'c' + abc$. Two neighbouring cells of a \omterm{def:b1:combinational-circuits:karnaugh}{Karnaugh map} differ in one variable, so their numbers of
1 inputs differ by one and their values of $f$ differ: no two 1 cells are neighbours, and no group is larger than one
cell.
\end{solution}

\begin{solution}{exo:b1:combinational-circuits:3}
The \omterm{def:b1:combinational-circuits:implicant}{essential} $cd'$ and $b'c'$ cover every 1 but the cells 5 and 7. Only $a'bd$ among the \omterm{def:b1:combinational-circuits:implicant}{prime implicants} covers both;
$a'c'd$ covers 5 and $a'bc$ covers 7, which needs two products. So $f = cd' + b'c' + a'bd$ is the only minimal sum; it
needs three products, since the two \omterm{def:b1:combinational-circuits:implicant}{essential} ones must be there and do not cover everything.
\end{solution}

\begin{solution}{exo:b1:combinational-circuits:4}
The code 11 is subtraction: $3 - 9 + 16 = 10 = 1010$, with $Z = 0$, $N = 1$, $C = 0$, $V = 1$. Unsigned, $C = 0$ means a
borrow: $3 < 9$ and the result has wrapped round. Signed, $9$ is $-7$ and $3 - (-7) = 10$ does not fit in 4 bits: $V =
1$, and the result read as $-6$ is wrong.
\end{solution}

\begin{solution}{exo:b1:combinational-circuits:5}
In the order $00, 01, 10, 11$, the middle columns $01$ and $10$ sit side by side but differ in two variables, $c$ and $d$:
they are not neighbours, and the two cells do not form a product. The \omterm{def:b1:combinational-circuits:canonical}{minterms} 1 and 2 stay apart: $f = a'b'c'd +
a'b'cd'$. The Gray order $00, 01, 11, 10$ makes the drawing tell the truth.
\end{solution}

\begin{solution}{exo:b1:combinational-circuits:6}
It computes $ab \oplus c(a + b)$. The two products are both 1 only when $a = b = 1$ and $c = 1$, and then XOR gives 0
where the carry is 1. On every other row at most one of them is 1, and XOR agrees with OR. The textbook carry $ab + c(a
\oplus b)$ has the same property, so there either gate works; with $a + b$ inside, only OR does.
\end{solution}

\begin{solution}{exo:b1:combinational-circuits:7}
\omcode{code/bachelor-1/27-combinational-circuits/exercises/mux.py}{10}{19}{\omterm{def:b1:combinational-circuits:mux}{Multiplexers}}
The 4-to-1 \omterm{def:b1:combinational-circuits:mux}{multiplexer} is two 2-to-1 \omterm{def:b1:combinational-circuits:mux}{multiplexers} on $s_0$ followed by one on $s_1$: 12 gates, delay 5. The test tries
the 64 inputs.
\end{solution}

\begin{solution}{exo:b1:combinational-circuits:8}
\omcode{code/bachelor-1/27-combinational-circuits/exercises/compare.py}{11}{17}{$a < b$ on four bits}
$a < b$ when, at some bit $i$, $a_i = 0$ and $b_i = 1$ while all higher bits are equal. One OR of four products: 20 gates,
delay 4 (XOR, NOT, AND, OR). The test tries the 256 pairs.
\end{solution}

\begin{solution}{exo:b1:combinational-circuits:9}
\omcode{code/bachelor-1/27-combinational-circuits/exercises/monotone.py}{7}{15}{Monotone functions, by \omterm{def:b1:backtracking:exhaustive}{brute force}}
The program finds 2, 3, 6 and 20 for $n = 0, 1, 2, 3$, against $2, 4, 16, 256$ functions in all. For $n = 1$: the
constants and $x$ ($x'$ is not monotone). For $n = 2$: the constants, $x$, $y$, $xy$ and $x + y$; any other function is
1 on some row and 0 on a row above it.
\end{solution}

\begin{solution}{exo:b1:combinational-circuits:10}
$x' = \text{NOR}(x, x)$ and $x + y = \text{NOR}(x, y)'$: NOR expresses NOT and OR, a complete set by De Morgan. For the
second part, every circuit of AND and OR gates computes a monotone function, by induction along a \omterm{def:b1:graph-traversals:topo}{topological order}:
the inputs are monotone, and AND and OR of monotone functions are monotone, since they are monotone in each argument.
NOT is not monotone, so no such circuit computes it.
\end{solution}

\begin{solution}{exo:b1:combinational-circuits:11}
In stage $i$, $p_i = a_i \oplus b_i$ and $a_i b_i$ have depth 1. Let $d_i$ be the depth of the carry $c_i$: $d_0 = 0$, and
$c_1 = a_0 b_0 + p_0 c_0$ has depth 3. For $i \geq 1$, $p_i c_i$ has depth $d_i + 1 \geq 4$ and the OR $d_i + 2$, so $d_i =
2i + 1$ for $i \geq 1$. The sums $s_i = p_i \oplus c_i$ have depth at most $d_i + 1 \leq 2n$, so the carry out, of depth
$2n + 1$, is the longest. In a lookahead block, $g_i$ and $p_i$ have depth 1, every AND of them (and of $c_0$) depth 2,
every carry, an OR, depth 3, and every sum depth 4.
\end{solution}

\begin{solution}{exo:b1:combinational-circuits:12}
\omcode{code/bachelor-1/27-combinational-circuits/exercises/multiplier.py}{12}{24}{An \omterm{def:b1:pointers-and-arrays:array}{array} multiplier}
The 16 partial products $a_i b_j$ cost 16 AND gates; three rows of four \omterm{def:b1:combinational-circuits:adder}{full adders} add them, row $j$ shifted $j$ places.
The circuit has 78 gates, including two that make a constant 0, and delay 18; the test compares it with $x \times y$ on
the 256 pairs.
\end{solution}

\begin{solution}{pb:b1:combinational-circuits:1}
\textbf{1.} On the rows $ab = 00, 01, 10, 11$: the sum is $0, 1, 1, 0$ and the carry $0, 0, 0, 1$.

\textbf{2.} Sum $a'b + ab'$, carry $ab$.

\textbf{3.} Read from left to right, the digits of \code{0b0011} are the values of $a$ on the rows $00, 01, 10, 11$, and
those of \code{0b0101} the values of $b$.

\textbf{4.} \code{&} works on the four rows in parallel, \code{\textasciitilde} inverts every bit (of an integer with
infinitely many leading bits, in Python), and \code{& 0b1111} keeps the four rows.

\textbf{5.} $2^{2^2} = 16$. One gate computes $\text{NAND}(a, a) = a'$, $\text{NAND}(b, b) = b'$ or $\text{NAND}(a, b)$:
three functions, the three circuits of one gate.

\textbf{6.} Each gate uses earlier signals only, so the circuit is a DAG, and its outputs are functions of its inputs.

\textbf{7.} $ab = \text{NAND}(n, n)$ with $n = \text{NAND}(a, b)$: two gates.

\textbf{8.} With $n = \text{NAND}(a, b)$: $\text{NAND}(\text{NAND}(a, n), \text{NAND}(b, n))$. On $00$: $n = 1$, both inner
gates give 1, the result 0; on $01$ and $10$: $n = 1$ and one inner gate gives 0, the result 1; on $11$: $n = 0$, both
inner gates give 1, the result 0.

\textbf{9.} Gate 0 is $n = \text{NAND}(a, b)$; gates 1 and 2 are $\text{NAND}(a, n)$ and $\text{NAND}(b, n)$; gate 3 is
$\text{NAND}(n, n) = ab$, the carry; gate 4 combines gates 1 and 2 into $a \oplus b$, the sum.

\textbf{10.} Both need $\text{NAND}(a, b)$: computing it once saves a gate, $2 + 4 - 1 = 5$.

\textbf{11.} Gate by gate: two XORs (4 each), two ANDs (2 each) and an OR (3), 15 NAND gates. Sharing, as in question 10,
brings a \omterm{def:b1:combinational-circuits:adder}{full adder} down to 9.

\textbf{12.} It grows: an AND becomes two levels, an OR two, an XOR three. NAND circuits are usually designed directly in
two levels: a sum of products is a NAND of NANDs, by De Morgan.

\textbf{13.} 4 bits: 20 gates and delay 9 rippling, 26 gates and delay 4 with lookahead; 8 bits: 40 and 17 against 52 and 6.

\textbf{14.} Rippling, each carry waits for the one before: two gates per bit. Lookahead computes each carry of a block
from the inputs by two levels of gates, whatever the block's position.

\textbf{15.} Between blocks the carries still ripple: the carry out of block $k$ is ready two delays after its carry in.
With $n / 4$ blocks the last sums come at $2(n / 4) + 2$: 6 for 8 bits, 34 for 64.

\textbf{16.} Five: the term $p_3 p_2 p_1 p_0 c_0$ of $c_4$. Gates with many inputs need more transistors in series and are
slower, so real designs limit them to about four inputs and use more levels.

\textbf{17.} $2^8 \times 2^8 \times 2 = 131\,072$: every possible input, so the check covers all behaviours of this
circuit.

\textbf{18.} $2^{129}$ inputs, beyond any computer: correctness for large $n$ needs a proof, by induction on the stages,
as in \cref{exo:b1:combinational-circuits:11}.

\textbf{19.} Gate $i$ may use any unordered pair, repetitions allowed, of the $i + 2$ signals before it:
$\binom{i + 3}{2}$ choices. The products $3$, $3 \times 6$, $\times 10$, $\times 15$, $\times 21$ are 3, 18, 180, 2700 and
56\,700, as printed.

\textbf{20.} $\text{NAND}(x, x) = x'$: without it, NOT would need a constant input.

\textbf{21.} A circuit of $j < k$ useful gates, padded with $k - j$ more, is among the circuits of $k$ gates; extra gates
do not change the signals already there.

\textbf{22.} Every circuit of at most four NAND gates appears among the 2700 tried (question 21), and none has both the
sum and the carry among its signals.

\textbf{23.} The search says nothing about them: it covers loop-free circuits of two-input NAND gates, which is the
claim.

\textbf{24.} \emph{The half adder needs exactly five NAND gates}: a circuit of five computes it (question 9), and the
\omterm{def:b1:backtracking:exhaustive}{exhaustive search} over all 2901 circuits of one to four two-input NAND gates finds none.
\end{solution}
